% Chapter body only — the template emits \chapter{A Taxonomy of Revivals in Quantum Walks}.
% Replace/extend the placeholder prose below. Keep \chapter{} out of this file.

\section{Why a taxonomy is needed}
The literature on revivals in quantum walks uses overlapping and sometimes
inconsistent terminology. ``Revival'', ``periodicity'', ``recurrence'', and
``return'' are all used, often for subtly different notions. This chapter
fixes vocabulary for the rest of the thesis.

\section{Revival types}
\subsection{Exact revivals}
\begin{definition}[Exact revivals]\label{def:exact-revival}
A walk $U$ acting on state $\psi_0$ exhibits an \emph{exact revival} at time $t > 0$ if $U^t \psi_0 = \psi_0$ up to a global phase.
\end{definition}
\begin{definition}[Exact revivals]\label{def:exact-revival}
A walk $U$ acting on state $\psi_0$ exhibits an \emph{exact revival} at time $t > 0$ if $U^t \psi_0 = \psi_0$ up to a global phase.
\end{definition}
\begin{definition}[Exact revivals]\label{def:exact-revival}
A walk $U$ acting on state $\psi_0$ exhibits an \emph{exact revival} at time $t > 0$ if $U^t \psi_0 = \psi_0$ up to a global phase.
\end{definition}

\subsection{Approximate revivals}
\begin{definition}[Approximate revivals]\label{def:approx-revival}
A walk exhibits an \emph{approximate revival} at time $t$ to tolerance $\epsilon$ if $\lVert U^t \psi_0 - \psi_0 \rVert < \epsilon$.
\end{definition}
\begin{definition}[Approximate revivals]\label{def:approx-revival}
A walk exhibits an \emph{approximate revival} at time $t$ to tolerance $\epsilon$ if $\lVert U^t \psi_0 - \psi_0 \rVert < \epsilon$.
\end{definition}
\begin{definition}[Approximate revivals]\label{def:approx-revival}
A walk exhibits an \emph{approximate revival} at time $t$ to tolerance $\epsilon$ if $\lVert U^t \psi_0 - \psi_0 \rVert < \epsilon$.
\end{definition}

\subsection{Fractional revivals}
\begin{definition}[Fractional revivals]\label{def:fractional-revival}
A walk exhibits a \emph{fractional revival} at time $t$ if $U^t \psi_0$ is supported on a strict subset of the initial support, with the non-zero amplitudes matching the initial ones up to phases.
\end{definition}
\begin{definition}[Fractional revivals]\label{def:fractional-revival}
A walk exhibits a \emph{fractional revival} at time $t$ if $U^t \psi_0$ is supported on a strict subset of the initial support, with the non-zero amplitudes matching the initial ones up to phases.
\end{definition}
\begin{definition}[Fractional revivals]\label{def:fractional-revival}
A walk exhibits a \emph{fractional revival} at time $t$ if $U^t \psi_0$ is supported on a strict subset of the initial support, with the non-zero amplitudes matching the initial ones up to phases.
\end{definition}

\subsection{Periodic revivals}
\begin{definition}[Periodic revivals]\label{def:periodic-revival}
A walk is \emph{periodic} on a subspace if exact revivals occur at a fixed period $T$ for every state in that subspace.
\end{definition}
\begin{definition}[Periodic revivals]\label{def:periodic-revival}
A walk is \emph{periodic} on a subspace if exact revivals occur at a fixed period $T$ for every state in that subspace.
\end{definition}
\begin{definition}[Periodic revivals]\label{def:periodic-revival}
A walk is \emph{periodic} on a subspace if exact revivals occur at a fixed period $T$ for every state in that subspace.
\end{definition}

\section{What is revived}
\subsection{State revivals}
\begin{definition}[State revivals]\label{def:state-revival}
The full state vector returns: $U^t \psi_0 = e^{i\theta} \psi_0$.
\end{definition}
\begin{definition}[State revivals]\label{def:state-revival}
The full state vector returns: $U^t \psi_0 = e^{i\theta} \psi_0$.
\end{definition}
\begin{definition}[State revivals]\label{def:state-revival}
The full state vector returns: $U^t \psi_0 = e^{i\theta} \psi_0$.
\end{definition}

\subsection{Probability revivals}
\begin{definition}[Probability revivals]\label{def:probability-revival}
Only the position distribution returns: $|(U^t \psi_0)(x)|^2 = |\psi_0(x)|^2$ for all $x$, with arbitrary phases on each amplitude.
\end{definition}
\begin{definition}[Probability revivals]\label{def:probability-revival}
Only the position distribution returns: $|(U^t \psi_0)(x)|^2 = |\psi_0(x)|^2$ for all $x$, with arbitrary phases on each amplitude.
\end{definition}
\begin{definition}[Probability revivals]\label{def:probability-revival}
Only the position distribution returns: $|(U^t \psi_0)(x)|^2 = |\psi_0(x)|^2$ for all $x$, with arbitrary phases on each amplitude.
\end{definition}

\subsection{Operator revivals}
\begin{definition}[Operator revivals]\label{def:operator-revival}
The evolution operator itself returns: $U^t = e^{i\theta} I$ on the relevant subspace.
\end{definition}
\begin{definition}[Operator revivals]\label{def:operator-revival}
The evolution operator itself returns: $U^t = e^{i\theta} I$ on the relevant subspace.
\end{definition}
\begin{definition}[Operator revivals]\label{def:operator-revival}
The evolution operator itself returns: $U^t = e^{i\theta} I$ on the relevant subspace.
\end{definition}

\section{The graph}
\subsection{Revivals on 1D cycles}
\begin{definition}[Revivals on 1D cycles]\label{def:revival-1d}
The graph is $C_N$, a cycle of length $N$, with a 2-state coin.
\end{definition}
\begin{definition}[Revivals on 1D cycles]\label{def:revival-1d}
The graph is $C_N$, a cycle of length $N$, with a 2-state coin.
\end{definition}
\begin{definition}[Revivals on 1D cycles]\label{def:revival-1d}
The graph is $C_N$, a cycle of length $N$, with a 2-state coin.
\end{definition}

\subsection{Revivals on 2D tori: separable vs.\ non-separable coins}
\begin{definition}[Revivals on 2D tori: separable vs.\ non-separable coins]\label{def:revival-2d}
The graph is $C_N \times C_M$. A separable coin factorises as $C = C_x \otimes C_y$; a non-separable coin does not.
\end{definition}
\begin{definition}[Revivals on 2D tori: separable vs.\ non-separable coins]\label{def:revival-2d}
The graph is $C_N \times C_M$. A separable coin factorises as $C = C_x \otimes C_y$; a non-separable coin does not.
\end{definition}
\begin{definition}[Revivals on 2D tori: separable vs.\ non-separable coins]\label{def:revival-2d}
The graph is $C_N \times C_M$. A separable coin factorises as $C = C_x \otimes C_y$; a non-separable coin does not.
\end{definition}

\section{A unified classification}
% TODO: big table summarising the taxonomy.
% Rows = revival type (exact / approximate / fractional / periodic)
% Columns = what is revived (state / probability / operator),
%           graph (cycle / torus), coin (separable / non-separable),
%           what the literature calls it.
%
% Example skeleton:
%
% \begin{table}[htbp]
%   \centering
%   \caption{A unified classification of quantum walk revivals.}
%   \label{tab:taxonomy}
%   \begin{tabular}{@{}llll@{}}
%     \toprule
%     Revival type & What revives & Graph & Literature term \\
%     \midrule
%     Exact    & State       & Cycle & ``full revival'' (Dukes) \\
%     ...      & ...         & ...   & ... \\
%     \bottomrule
%   \end{tabular}
% \end{table}
